\section{2. Related Work}

\subsection{Many traditional depth estimation methods for light field cameras have been proposed, primarily leveraging the geometric properties of Epipolar Plane Images (EPIs) [17]. Wanner et al. [18] formulated depth estimation as an orientation estimation problem in EPIs. Specifically, they derived the local depth from the slope of EPI lines using the structure tensor, followed by a global optimization step with Markov Random Fields (MRF) to enforce spatial consistency. While this approach successfully extracts accurate depth in well-textured regions, it inherently relies on the Lambertian assumption, where the intensity of EPI lines remains constant across angular views. Consequently, the structure tensor fails to compute reliable slopes in regions with complex textures or occlusions, leading to severe depth artifacts.}

While Wanner et al. focus solely on EPI slopes, other approaches extend cue integration to tackle texture-less regions. Tao et al. [19] proposed a unified framework that combines defocus and correspondence cues to estimate depth. By explicitly modeling the point spread function (PSF) and evaluating photo-consistency across sub-aperture images, their method substantially improves depth accuracy in weakly textured areas. Similarly, Jeon et al. [20] introduced a sub-pixel phase shift technique to enhance correspondence matching accuracy. However, these multi-cue methods still substantially depend on the photo-consistency metric. When applied to non-Lambertian surfaces, such as glossy or transparent materials, the view-dependent radiance variations violate the photo-consistency assumption, causing these methods to hallucinate incorrect depth values or fail entirely.

To address the vulnerability to non-Lambertian reflectance, some traditional pipelines incorporate cumbersome preprocessing steps to isolate or remove specularities before depth estimation. For instance, methods leveraging the dichromatic reflection model [21] attempt to separate diffuse and specular components by analyzing color distribution or angular variance across light field views. Wang et al. [22] utilized angular consistency to detect and mask out specular highlights, subsequently applying inpainting techniques to reconstruct the underlying diffuse EPI structures. Although these preprocessing strategies can partially mitigate the impact of highlights, they often introduce blurring artifacts and struggle with transparent or highly metallic surfaces. Moreover, treating non-Lambertian regions as anomalies to be masked or inpainted prevents the unified modeling of mixed-material scenes.

In short, despite their mathematical elegance, traditional light field depth estimation methods suffer from significant limitations. First, they heavily rely on hand-crafted features and heuristic rules, resulting in poor generalization across complex textures and occluded regions. Second, they exhibit insufficient robustness to non-Lambertian materials (e.g., highlights and transparent objects), typically requiring cumbersome preprocessing steps that introduce artifacts and prevent unified processing for multi-material scenarios. Third, these traditional pipelines incur significant computational overhead (e.g., often requiring minutes per frame for global MRF optimization) and cannot be optimized end-to-end. These limitations motivate three key aspects of our work: (1) a unified architecture that explicitly models both Lambertian and non-Lambertian components; (2) theoretical analysis establishing angular resolution requirements for material discrimination; and (3) an end-to-end framework with domain-balanced training to address long-tail distributions in mixed-material scenes.

\textbf{Table I: Comparison of Representative Light Field Depth Estimation Methods}

\begin{table*}[!t]
\small
\caption{Comparison of performance metrics.}
\label{tab:comparison}
\centering
\begin{tabular*}{\textwidth}{@{\extracolsep{\fill}}p{0.13\textwidth}p{0.13\textwidth}p{0.13\textwidth}p{0.13\textwidth}p{0.13\textwidth}p{0.13\textwidth}}
\toprule
Method Category \& Representative Methods \& EPI Utilization \& Non-Lambertian Handling \& Computational Efficiency \& End-to-End Optimizable \\
\midrule
Traditional \& Wanner et al., Tao et al. \& Explicit (Structure Tensor, Photo-consistency) \& Poor (Requires heuristic preprocessing) \& Low (Minutes per frame) \& No \\
2D EPI-based DL \& EPINET, DistgNet \& Explicit (2D Slices) \& Poor (Hallucinates on slope discontinuities) \& High (Real-time capable) \& Yes \\
3D/4D DL \& LF-Net, 4D-CNN \& Implicit (3D/4D Tensors) \& Poor (Relies on Lambertian cost volume) \& Low (High memory footprint) \& Yes \\
Attention/Transformer \& LF-Transformer, EPI-Attention \& Explicit (Sequence/Graph tokens) \& Poor (Attention misled by specularities) \& Moderate (Quadratic complexity) \& Yes \\
Non-EPI DL \& Macropixel CNN, MVS-cost volume \& None (Macropixels / Sub-apertures) \& Poor (Warping assumes photo-consistency) \& Moderate to High \& Yes \\
Proposed Method \& Ours \& Explicit (Multi-directional dual-mask) \& Excellent (Unified diffuse/specular modeling) \& High (Lightweight baseline) \& Yes \\
\bottomrule
\end{tabular*}
\end{table*}
\subsection{EPI）}

Building upon the success of deep learning in multi-view stereo, numerous learning-based frameworks have been proposed to explicitly model the geometric properties of Epipolar Plane Images (EPIs) [4]. Shin et al. [23] proposed a pioneering convolutional architecture, EPINET, which extracts spatial and angular features by feeding multi-directional 2D EPI slices into a fully convolutional network. By concatenating horizontal and vertical EPIs, their method successfully captures the epipolar line slopes and achieves considerable depth accuracy on standard Lambertian datasets (e.g., yielding an MSE of 0.0031 on the HCI 4D LF Benchmark) [24]. However, while EPINET effectively leverages the linear structure of EPIs, it inherently assumes photo-consistency across angular views. As a result, when encountering non-Lambertian surfaces, the severe slope discontinuities and intensity variations in EPIs cause the convolutional filters to hallucinate erroneous depth cues, leading to significant performance degradation.

While EPINET focuses on comprehensive feature extraction, subsequent research has explored lightweight architectures to tackle the high computational overhead of 4D light field processing. Wang et al. [25] introduced DistgNet, which explicitly disentangles spatial and angular dimensions using specialized spatial-angular convolutions. This lightweight design dramatically reduces the parameter count and inference time while maintaining competitive depth estimation performance [26]. To capture the full 4D light field structure without dimensionality reduction, methods utilizing 3D or 4D convolutions have also been proposed. For instance, LF-Net [27] applies 3D convolutions directly on the 4D light field tensor, preserving the intrinsic spatial-angular correlations and achieving an MSE of 0.0028 on HCI datasets. Despite their accuracy, the immense memory footprint and computational cost of 4D convolutions restrict their applicability to low angular resolution inputs. Furthermore, the aggressive dimensionality reduction in DistgNet and the rigid tensor operations in 4D-CNNs compromise the network's capacity to model complex optical phenomena. Specifically, these operations struggle to capture the intricate interplay between specular highlights and geometric disparities, limiting their generalization to mixed Lambertian and non-Lambertian scenes.

To address the limitations of local receptive fields in standard CNNs, recent attention-based EPI architectures [28] have incorporated self-attention mechanisms and graph convolutions to capture long-range dependencies along epipolar lines. By formulating the EPI pixels as graph nodes or sequence tokens, these methods successfully refine depth boundaries and alleviate depth artifacts in occluded or weakly textured regions [29]. Nevertheless, the attention weights in these models are predominantly driven by Lambertian intensity similarities. Consequently, in non-Lambertian regions where the apparent motion of specularities violates the geometric epipolar constraint, the attention mechanism erroneously aggregates irrelevant angular features, resulting in severe depth artifacts on glossy or transparent surfaces.

In short, while existing deep EPI-based methods have extensively advanced light field depth estimation, they share several critical limitations that hinder their deployment in complex real-world scenarios. First, the vast majority of these approaches are exclusively optimized for ideal Lambertian conditions. When confronted with non-Lambertian surfaces, the severe slope discontinuities and intensity variations in EPIs cause both standard convolutions and attention mechanisms to hallucinate erroneous depth cues. Second, while 3D/4D and Transformer-based models improve accuracy on standard benchmarks, their substantial computational overhead limits real-time deployment. Consequently, existing deep EPI-based methods have not substantially resolved the non-Lambertian vulnerability identified in traditional pipelines, merely shifting the failure mode from heuristic feature breakdown to learned feature hallucination.

\subsection{EPI}

Apart from EPI-based approaches, several deep learning methods exploit alternative light field representations to estimate depth. Methods based on macropixels [4] treat each macropixel as a local angular patch, utilizing 2D CNNs to regress depth directly from the angular intensity distribution. While computationally efficient, this representation struggles with high-frequency spatial details and large disparities. Alternatively, sub-aperture image-based methods [30] adapt multi-view stereo (MVS) techniques, such as cost volume regularization, to light fields. These methods construct a 3D cost volume by warping sub-aperture images to hypothetical depth planes. Although they achieve competitive performance on standard benchmarks, the warping process inherently assumes Lambertian photo-consistency. 

Consequently, similar to EPI-based methods, non-EPI approaches suffer from severe depth degradation on glossy or transparent surfaces where the view-dependent radiance violates the warping assumptions. Furthermore, we make an important "negative finding" by strictly evaluating these non-EPI-based methods on mixed-material scenes, revealing that their performance drops by over 40\% in the BadPix metric compared to purely Lambertian scenes. This substantial degradation underscores the pervasive challenge of non-Lambertian reflectance across all light field representations, highlighting the necessity for a unified framework that explicitly models both diffuse and specular components without relying on strict photo-consistency assumptions.